\documentclass[a4paper,10pt]{article}

%----------------------------------------------------------------------------------------
%	PACKAGES AND THEMES
%----------------------------------------------------------------------------------------
\usepackage[top=0.5in, bottom=0.5in, left=0.5in, right=0.5in]{geometry}
\usepackage[T1]{fontenc}
\usepackage{enumitem}
\usepackage{hyperref}
\usepackage{xcolor}
\usepackage{titlesec}
\usepackage{fontawesome5}
\usepackage{tabularx}

% --- COLOR: AQUA BLUE ---
\definecolor{primary}{RGB}{0, 160, 210} 

% --- LINKS: NO BOXES, COLORED TEXT ---
\hypersetup{
    colorlinks=true,
    urlcolor=primary,
    linkcolor=primary
}

% Formatting Sections
\titleformat{\section}
  {\Large\bfseries\color{primary}\uppercase} 
  {} 
  {0em} 
  {}[\titlerule] 

\titlespacing{\section}{0pt}{10pt}{5pt}

% --- CUSTOM COMMANDS FOR ROLES ---
% Format: Role | Place | Date (All Bold)
%         Advisor (Below)
\newcommand{\role}[4]{
    \noindent \textbf{#1 | #3 | #2} \\
    \textit{#4}
}

% Publication Command
\newcommand{\pub}[3]{
    \noindent \textbf{#1} \\
    #2 \\
    \textit{#3}
}

% General Setup
\pagestyle{empty} 
\setlength{\parindent}{0pt}
\setlist[itemize]{leftmargin=*, noitemsep, topsep=0pt}

%----------------------------------------------------------------------------------------
%	DOCUMENT START
%----------------------------------------------------------------------------------------
\begin{document}

%--- HEADER ---
\begin{center}
    {\Huge \textbf{\color{primary} Omar Muhammad}} \\ \vspace{4pt}
 \ \href{mailto:omarmuhammad@iisc.ac.in}{omarmuhammad@iisc.ac.in} \ | \ \href{https://omar826.github.io/}{Personal Website}
\end{center}

%--- PROFILE ---

%--- EDUCATION ---
\section*{Education}
% Format: Degree | Date | Institute | CGPA (All Bold)
\noindent \textbf{B-Tech in Mathematics and Computing | 2023 -- 2027 | Indian Institute of Science, Bengaluru | CGPA: 9.5/10.0}

\vspace{3pt}
\noindent \textbf{Grade 12 (CBSE) | 2022 -- 2023 | Saraswathi Vidyanikethan, Ernakulam | 98.4\% (School Top)}

\vspace{3pt}
\noindent \textbf{Grade 10 (CBSE) | 2020 -- 2021 | Rajagiri Public School, Ernakulam | 98\%}

%--- PUBLICATIONS ---
\section*{Publications}
\pub{Verification Modulo Tested Library Contracts}
{Abhishek Uppar, Omar Muhammad, Sumanth Prabhu S, Deepak D'Souza, P. Madhusudan, Adithya Murali.}
{Accepted to ACM SIGPLAN Conference on Programming Language Design and Implementation (PLDI 2026) and appeared in PACMPL. \href{https://dl.acm.org/doi/10.1145/3808305}{[Paper]} \href{https://github.com/omar826/Dualis-VMTLC}{[Code]}} 

\pub{Have a thing? Reasoning around recursion with dynamic typing in grounded arithmetic}
{Bryan Ford, Elliot Bobrow, Stefan Milenković, Omar Muhammad, Dimitrios Alexopoulos, Yusuf Demir, Ananthajit Srikant}
{Under review at the ACM SIGPLAN Conference on Principles of Programming Languages.
\href{https://arxiv.org/abs/2510.25369}{[Paper]}
}

\pub{Mask2Cause: Causal Discovery via Adjacency Constrained Causal Attention}
{Omar Muhammad, Pasupuleti Dhruv Shivkant, Deepak N. Subramani}
{Under review at the conference on Neural Information Processing Systems (NeurIPS 2026). \href{https://arxiv.org/abs/2605.07280}{[Paper]} \href{https://github.com/omar826/Mask2Cause}{[Code]}}

%--- RESEARCH EXPERIENCE ---
\section*{Research Experience}

\role{Summer Research Fellow, Grounded Deduction}{May 2026 -- Present}{EPFL, Switzerland (Onsite for May – July 2026; Remote after)}{Adviser: Prof. Bryan Ford, DEDIS Lab}
\begin{itemize}
    \item Developing Grounded Arithmetic, a non-classical logic that admits arbitrary recursive definitions without prior termination proofs, and mechanising its metatheory in Isabelle. \href{https://omar826.github.io/grounded.pdf}{[Slides]}
    \item Proved truth preservation and primitive recursiveness for the inference rules of Basic Grounded Arithmetic (BGA). Contributed to the design of Propositional Grounded Arithmetic (PGA).
    \item Built a quantified extension of Grounded Arithmetic as a usable object logic in Isabelle/Pure, and established the consistency of BGA within it.
    \item Designing automated proof-search tactics for the Isabelle/Pure development, which inherits none of the automation built for classical and intuitionistic logics.
\end{itemize}

\vspace{5pt}
\role{Student Researcher, Verification Modulo Tested Library Contracts}{May 2025 -- Apr 2026}{IISc}{Advisers: Prof. D. D'Souza (IISc), Prof. P. Madhusudan (UIUC), Prof. A. Murali (UW-Madison)}
\begin{itemize}
    \item Co-developed VMTLC, a hybrid framework that verifies client programs by synthesizing library specifications that are formally proven against the client but validated via testing against the library.
    \item Designed the Learner component within a Teacher-Learner CEGIS loop, developing two distinct LLM-integrated architectures (a pure LLM learner and an LLM-seeded ICE learner) to leverage semantic intuition, synthesizing valid contracts where traditional enumerative baselines could not.
    \item Designed and built the core LLM pipeline responsible for generating contracts from logical rules and integrated it into a feedback loop with formal verifiers and testing engines to ensure correctness.
\end{itemize}

\vspace{5pt}
\role{Research Intern, Enhancing Formal Verification with LLMs}{Aug 2025 -- Oct 2025}{University of Virginia, USA (Remote)}{Adviser: Prof. Wenxi Wang, The University of Virginia}
\begin{itemize}
    \item Developed ACPRO, a framework that uses a multi-agent LLM pipeline to transform low-level Z3 proofs into interpretable, high-level Verus verification steps.
    \item This work also produced the VSVerus dataset, a new resource for studying fine-grained LLM reasoning in formal verification.
\end{itemize}

\vspace{5pt}
\role{Summer Research Fellow, Formalization of Group Theoretic Algorithms}{May 2025 -- July 2025}{Krea University, India (Onsite)}{Advisers: Prof. Rishi Vyas (KREA), Prof. T.V.H Prathamesh (KREA)}
\begin{itemize}
    \item Formalized proofs for decision problems in geometric group theory (including the word, conjugacy, and subgroup membership problems) using the Lean 4 proof assistant.
    \item Key contributions include formalizing two definitions of the Dehn function and proving their equivalence and verifying the algorithm for the conjugacy problem in free groups.
    \item Current work focuses on the computability of the Dehn function for hyperbolic groups.
\end{itemize}

\vspace{5pt}
\noindent
\textbf{Student Researcher, Game Theory and Automata | Jan 2026 -- Present | IISc}
\begin{itemize}
    \item Investigating supergames played by finite automata under finite costs for complexity to analyze the sequential collapse of equilibrium states.
    \item Deriving critical phase change thresholds that characterize these equilibrium dynamics.
\end{itemize}

\vspace{5pt}

\role{Student Researcher, Causal Discovery in Time Series}{Aug 2025 -- Present}{IISc}{Adviser: Prof. Deepak N. Subramani, Computational and Data Sciences, IISc}
\begin{itemize}
    \item Proposed and developed Mask2Cause, an end-to-end Transformer that recovers the causal graph during the forecasting forward pass via a learnable sparse adjacency injected as a log-barrier into masked attention over inverted variable embeddings.
    \item Formulated causal discovery under volatility spillover, where one variable governs the conditional variance rather than the mean of another. Motivated the heteroscedastic objective through the theory of directed information graphs and showed neural Granger causality to be its additive-noise special case.
    \item Built Mixed Physics, the first benchmark separating mean-parents from variance-parents, and the first model to recover variance-mediated edges explicitly. Attained state-of-the-art accuracy at reduced parameter count.
\end{itemize}

%--- TEACHING ---
\section*{Teaching}
\role{Teaching Assistant, Program Analysis and Verification}{Aug 2026 -- Present}{IISc}{Instructors: Prof. Deepak D'Souza, Prof. K. V. Raghavan}
\begin{itemize}
    \item Designed the course project and mentored student teams through it. Graded assignments, conducted vivas, and held regular office hours for doubt clearing.
\end{itemize}

%--- PROJECTS ---
\section*{Projects}

\role{Independent Study, Logic and Automated Theorem Proving}{Aug 2026 -- Present}{IISc}{Adviser: Prof. Deepak D'Souza}
\begin{itemize}
    \item Studying structural proof theory (sequent calculi, cut-elimination and the properties a calculus needs for proof search to be feasible), provability and self-reference (G\"odel's incompleteness theorems, the derivability conditions, L\"ob's and Tarski's theorems), decidable fragments of logic, and automated reasoning (unification, resolution, superposition, tableaux and decision procedures for decidable theories). \href{https://omar826.github.io/ISP_notes.pdf}{[Notes: Completeness of First-Order Logic]}
\end{itemize}
\vspace{3pt}

\noindent \textbf{Reading Project, Threshold for Random k-SAT | Jan 2026 -- April 2026}
\begin{itemize}
    \item Explored the satisfiability threshold for random k-SAT as part of the ``Techniques in Discrete Probability'' course under the guidance of Prof. Riddhipratim Basu, ICTS-TIFR.
    \item Analyzed the foundational work by Dimitris Achlioptas and Yuval Peres, on the application of weighted moment methods to bound the threshold. \href{https://omar826.github.io/report_ksat.pdf}{[Report]} \href{https://omar826.github.io/presentation_ksat.pdf}{[Presentation]}
\end{itemize}

\vspace{3pt}
\noindent \textbf{Reading Project, PCP Theorem | Aug 2025 -- Nov 2025}
\begin{itemize}
    \item Conducted an in-depth study of the PCP theorem, analyzing its formulations in terms of both probabilistically checkable verifiers and the hardness of approximation under the guidance of Prof. Anand Louis, Department of Computer Science and Automation, IISc.
    \item Analyzed Irit Dinur's proof of the theorem using gap amplification. \href{https://omar826.github.io/PCP_report.pdf}{[Report]}
\end{itemize}

\vspace{3pt}
\noindent \textbf{Formalization Project, Rubik's Cube Solvability Conditions | Feb 2025 -- May 2025}
\begin{itemize}
    \item Modeled the nxnxn Rubik's Cube as a group, formalized proof of validity of solvability conditions of the $4\times4\times4$ Rubik's Cube and formulated the solvability conditions for the nxnxn cube in Lean 4. \href{https://github.com/omar826/N_Rubiks_Cube}{[Github Repo]}
\end{itemize}

\vspace{3pt}
\noindent \textbf{Reading Project, Undecidability of First-Order Logic | Feb 2025 -- April 2025}
\begin{itemize}
    \item Studied the foundational proof by Alonzo Church, and its implications for computability and formal verification systems.
\end{itemize}

\vspace{3pt}
\noindent \textbf{Research Project, Image Denoising Algorithms | March 2025 -- April 2025}
\begin{itemize}
    \item Implemented and evaluated various frameworks for image restoration without clean data.
    \item Designed and conducted experiments to analyze the hyperparameter sensitivity of Zero-Shot Noise2Noise implementation. \href{https://omar826.github.io/noise_report.pdf}{[Report]} \href{https://omar826.github.io/noise_presentation.pdf}{[Presentation]}
\end{itemize}

\vspace{3pt}
\noindent \textbf{Reading project, Differential Topology and Knot Theory | May 2024 -- July 2024}
\begin{itemize}
    \item Explored topics in Differential Topology and Knot Theory under the guidance of Prof. Subhojoy Gupta, Department of Mathematics, Indian Institute of Science.
\end{itemize}

\vspace{3pt}
\noindent \textbf{Developer DataBased (CS club at IISc) | Oct 2023 -- Dec 2023}
\begin{itemize}
    \item Utilized Generative AI to develop components for an automated Refute Problem Generation system.
\end{itemize}

\section*{Software}
\role{Intern}{Dec 2024 -- June 2025}{National Payments Corporation of India (NPCI)}{Adviser: Prof. Viraj Kumar, Divecha Centre for Climate Change, IISc}
\begin{itemize}
    \item Developed an internal tool to enable NPCI employees to query databases and generate data visualizations via natural language prompts.
    \item Implemented a pipeline that translates natural language queries to SQL queries, fetches data, determines correct visualization types, and renders visual output.
\end{itemize}

\vspace{5pt}
\role{Intern}{July 2024 -- Jan 2025}{Applied Research Capability Network (ARC-Net), IISc Campus}{}
\begin{itemize}
    \item Fine-tuned LLMs with Tool Integrated reasoning for solving math and physics problems.
    \item Used proof assistants and symbolic solvers to verify the chain of thought of the models.
\end{itemize}

%--- ACHIEVEMENTS ---
\section*{Achievements and Recognitions}
\begin{itemize}
    \item Selected for MITACS Globalink Research Internship 2026 (Offers: Undergraduate Summer Research Program [UGSRP] at University of Toronto; Fields Undergraduate Summer Research Program [FUSRP] at University of Waterloo)
    \item JEE Advanced Examination: 646th rank
    \item JEE Mains Examination: 99.95 percentile, 654th rank
    \item KVPY (Kishore Vaigyanik Protsahan Yojana) Scholar
    \item NTSE (National Talent Search Examination) Scholar
\end{itemize}

%--- TECHNICAL SKILLS ---
\section*{Technical Skills}
\begin{itemize}
    \item \textbf{Verification \& Logic:} Isabelle, Lean 4, Z3, SMT Solvers, Alloy, Spin, Rodin, VCC, JPF, Verus
    \item \textbf{Mathematics:} Statistics, Probability, Topology, Abstract Algebra, Real Analysis, Linear Algebra, Combinatorics, Graph Theory, Number Theory, Multivariable Calculus, Stochastic Processes
    \item \textbf{Theoretical Computer Science:} Automata Theory, Game Theory \& Mechanism Design, Program Analysis \& Verification, Numerical Methods, Probabilistic Methods, Information Theory, Discrete Fourier Analysis, Convex Optimization
    \item \textbf{Machine Learning:} PyTorch, Scikit-learn, Transformers, LLM Fine-Tuning, Causal Discovery
    \item \textbf{Languages \& Tools:} Python, C, SML, SQL

\end{itemize}

%--- EXTRACURRICULAR ---
\section*{Extracurricular Activities}
\begin{itemize}
    \item Mentored underprivileged students in STEM concepts through the Notebook Drive Mentorship Outreach Program, IISc.
    \item Active Member of DataBased (Computer Science Club at IISc).
    \item Designed competitive programming problems for the RevCoding Contest as a Volunteer for Pravega 2025 (IISc UG Fest).
    \item Presented on 'Online Decision-Making Algorithms' at Algorithm Festival 2024, IISc.
    \item Enthusiastic participant in MUN (Model United Nations) conferences and debate competitions.
\end{itemize}

\section*{Languages}
\begin{itemize}
    \item \textbf{Fluent:} English, Malayalam
    \item \textbf{Intermediate:} Hindi
\end{itemize}

\end{document}